\documentclass[12pt,letterpaper]{article}

% Packages
\usepackage[margin=1in]{geometry}
\usepackage{graphicx}
\usepackage{amsmath}
\usepackage{amssymb}
\usepackage{cite}
\usepackage{hyperref}
\usepackage{caption}
\usepackage{subcaption}
\usepackage{float}
\usepackage{booktabs}

% Title and author
\title{\textbf{Final Project Report: Beam Training with Analog True-Time-Delay Arrays}}
\author{Jake Sherry - 805745388}
\date{September 4, 2026}

\begin{document}

\maketitle

\section{Introduction}

This report presents the implementation and evaluation of a single-shot beam training algorithm for wideband millimeter-wave (mmWave) orthogonal frequency division multiplexing (OFDM) systems using true-time-delay (TTD) arrays. MmWave communications are essential for achieving high data rates in fifth-generation and beyond wireless systems due to abundant available spectrum. However, these systems face severe propagation loss and require beamforming with large antenna arrays at both the base station (BS) and user equipment (UE) to maintain reliable links~\cite{boljanovic2020design}.

Conventional phased-array systems based on phase shifters require exhaustive beam sweeping to identify dominant propagation directions, resulting in significant training overhead. With such arrays, all frequency components are steered to or combined from the same direction because identical phase shifts are applied across all signal frequencies. As array sizes increase and beamwidths narrow, this training overhead becomes prohibitive.

TTD arrays offer a fundamentally different approach. Unlike phase shifters, TTD elements introduce frequency-dependent delays that enable different frequencies to probe different angular directions simultaneously~\cite{boljanovic2020design}. This frequency-diversity property allows UE equipped with a TTD array to determine the dominant propagation direction using a single training symbol.

In this project, I attempted to implement and evaluate the TTD-based beam training system proposed by Boljanovic et al.~\cite{boljanovic2020design}. Based on their work, I designed frequency-dependent beam training codebooks for analog TTD arrays and implemented a dictionary-based angle estimation algorithm. The system operates at 60~GHz with 2~GHz bandwidth and uses 128 base station antennas and 16 user equipment antennas. To measure the perofrmance, I evaluated the root mean square error (RMSE) of angle of arrival estimation under various signal-to-noise ratios and hardware impairment conditions, including gain, phase, and delay errors.

The results in the paper demonstrate that the TTD-based approach enables efficient beam training with diversity factors of $R = 1, 2, 4$, achieving accurate angle estimation. However, I observed several discrepancies compared to the reference paper, including an elevated noise floor in my SNR results and a stronger sensitivity to phase errors relative to gain errors, suggesting implementation differences despite my best efforts to adhere as closely as possible to the mathematical formulations in~\cite{boljanovic2020design}.

\section{System Model}

\subsection{OFDM-Based Wideband Communication}

As outlined in the project spec, I considered a downlink beam training scenario between a BS and UE using OFDM with cyclic prefix. The system parameters follow the specifications in~\cite{boljanovic2020design} and are summarized in Table~\ref{tab:parameters}.

\begin{table}[h]
\centering
\caption{System Parameters}
\label{tab:parameters}
\begin{tabular}{@{}lcc@{}}
\toprule
\textbf{Parameter} & \textbf{Symbol} & \textbf{Value} \\
\midrule
Carrier Frequency & $f_c$ & 60 GHz \\
Bandwidth & BW & 2 GHz \\
OFDM Subcarriers & $M_{\text{tot}}$ & 4096 \\
BS Antennas & $N_T$ & 128 \\
UE Antennas & $N_R$ & 16 \\
Training Directions & $D$ & 32 \\
Dictionary Size & $Q$ & 1024 \\
Channel Clusters & $L$ & 3 \\
Power Ratio & $\sigma_1^2/\sigma_2^2$ & 10 dB \\
Monte Carlo Trials & $N_{\text{trials}}$ & 1000 \\
\bottomrule
\end{tabular}
\end{table}

\noindent
The BS is equipped with a digital array and uses a fixed precoder $\mathbf{v} = \mathbf{a}_T(\hat{\theta}^{(T)})$ designed according to a known angle of departure, where $\mathbf{a}_T(\theta)$ is the transmit array response vector. The UE employs a TTD array with delay elements implemented in the baseband domain, enabling frequency-dependent beam patterns as described in~\cite{boljanovic2020design}.

\subsection{Channel Model}

I implemented a wideband geometric channel model with $L = 3$ multipath clusters, where the dominant path is 10~dB stronger than each of the other two paths. For a given sub-band $k$, the channel matrix is expressed as:
\begin{equation}
\mathbf{H}[k] = \sum_{l=1}^{L} G_l[k] \mathbf{a}_R(\theta_l^{(R)}) \mathbf{a}_T^H(\theta_l^{(T)})
\end{equation}
where $\theta_l^{(R)}$ and $\theta_l^{(T)}$ are the angle of arrival (AoA) and angle of departure (AoD) of the $l$-th cluster, and $G_l[k]$ represents the complex channel gain. The array response vectors for a half-wavelength spaced uniform linear array are:
\begin{equation}
[\mathbf{a}_R(\theta)]_n = N_R^{-1/2} \exp[j(n-1)\pi\sin(\theta)], \quad n = 1, \ldots, N_R
\end{equation}

\noindent
The complex gains $G_l[k]$ are modeled as independent zero-mean complex Gaussian random variables with variance $\sigma_l^2$, where $\sigma_1^2 \geq \sigma_2^2 \geq \sigma_3^2$ and $\sigma_1^2/\sigma_2^2 = 10$~dB.

\subsection{Received Signal with TTD Array}

The BS transmits a power-normalized training pilot on $M$ selected subcarriers: $X[m] = M^{-1/2}$ for $m \in \mathcal{M}$. At the UE, the received signal at the $m$-th subcarrier is:
\begin{equation}
Y[m] = M^{-1/2} \mathbf{w}^H[m] \mathbf{H}[k] \mathbf{v} + \mathbf{w}^H[m] \mathbf{n}[m]
\end{equation}
where $\mathbf{n}[m] \sim \mathcal{CN}(0, \sigma_N^2 \mathbf{I}_{N_R})$ is additive white Gaussian noise, and $\mathbf{w}[m]$ is the frequency-dependent TTD antenna weight vector.
\\
\\
\textbf{Note:} Based on my understanding, the reference paper~\cite{boljanovic2020design} does not explicitly specify how noise power should be scaled relative to the number of training subcarriers $M$. In my implementation, I include the factor $M$ in the noise power calculation (i.e., the noise variance is scaled by $M$) based on the assumption that this accounts for the energy aggregation across $M$ subcarriers.

\section{TTD Codebook Design}

\subsection{Frequency-Dependent Beamforming}

The key advantage of TTD arrays lies in the ability to synthesize frequency-dependent beams. The $n$-th element of the TTD weight vector at subcarrier $m$ is:
\begin{equation}
[\mathbf{w}[m]]_n = \exp[j(2\pi f_m \tau_n + \phi_n)]
\end{equation}
where $f_m = f_c - \text{BW}/2 + (m-1)\text{BW}/(M_{\text{tot}}-1)$ is the subcarrier frequency, $\tau_n$ is the delay tap, and $\phi_n$ is the phase tap for the $n$-th antenna element.

\subsection{Diversity Factor and Subcarrier Selection}

To introduce frequency diversity, the bandwidth is divided into $R$ distinct sub-bands. Within each sub-band, there are $D$ uniformly spaced subcarriers associated with $D$ fixed directional beams. For a given direction $d$, the set of $R$ subcarriers mapped to that direction is:
\begin{equation}
\mathcal{M}_d = \left\{ m \mid m = 1 + (d-1)\lfloor M_{\text{tot}}/(DR)\rfloor + (r-1)M_{\text{tot}}/R, \; r = 1, \ldots, R \right\}
\end{equation}
The parameter $R$ is called the diversity factor, representing the number of independent frequency samples obtained for each direction.

\subsection{Delay and Phase Tap Design}

To achieve the desired frequency-dependent beam pattern with $D = 32$ directions and diversity factor $R$, uniformly spaced delay and phase taps are designed as:
\begin{align}
\tau_n &= (n-1)R/\text{BW} \\
\phi_n &= (n-1)[\text{sgn}(\psi)\pi - \psi]
\end{align}
where $\psi = \text{mod}(2\pi R(f_c - \text{BW}/2)/\text{BW} + \pi, 2\pi) - \pi$.
\\ \\
This design ensures that each subcarrier in $\mathcal{M}_d$ produces a beam pointing in direction $d$, with the beam pattern given by a discrete Fourier transform (DFT) beam:
\begin{equation}
[\mathbf{f}_d]_n = \exp[-j2\pi(n-1)(d-1-D/2)/D], \quad d \leq D
\end{equation}

\section{Angle Estimation Algorithm}

\subsection{Power-Based Super-Resolution Method}

The algorithm estimates the AoA by analyzing the received power across all $D$ directions. For direction $d$, the received signal power from the $R$ subcarriers in $\mathcal{M}_d$ provides independent samples of the channel. The estimate of the power in direction $d$ is:
\begin{equation}
\hat{p}_d = \frac{1}{R} \sum_{m \in \mathcal{M}_d} |Y[m]|^2
\end{equation}

\subsection{Dictionary-Based Estimation}

The expected power vector $\mathbf{p} = [p_1, \ldots, p_D]^T$ can be related to a dictionary matrix $\mathbf{B}$ via:
\begin{equation}
\mathbf{p} = \mathbf{B}\mathbf{g} + N_R\sigma_N^2 \mathbf{1}
\end{equation}
where $\mathbf{B} \in \mathbb{R}^{D \times Q}$ is constructed using a grid of $Q = 1024$ uniformly spaced candidate angles $\{\xi_q\}_{q=1}^Q$:
\begin{equation}
[\mathbf{B}]_{d,q} = |\mathbf{f}_d^H \mathbf{a}_R(\xi_q)|^2
\end{equation}

\noindent
The vector $\mathbf{g}$ is sparse with significant energy only at the index corresponding to the dominant cluster's AoA. The AoA estimate is obtained by finding the dictionary column with highest correlation to the measured power vector:
\begin{equation}
\hat{\theta}^{(R)} = \xi_{q^*}, \quad \text{where} \quad q^* = \arg\max_q \frac{\hat{\mathbf{p}}^T [\mathbf{B}]_{:,q}}{\|[\mathbf{B}]_{:,q}\|}
\end{equation}

\section{Hardware Impairments}

Real TTD arrays suffer from hardware impairments that degrade beam training performance. I considered three types of errors as modeled in~\cite{boljanovic2020design}:

\subsection{Gain Error}
Magnitude mismatches across antenna elements are modeled with log-normal distribution:
\begin{equation}
10\log_{10}(\alpha_n) \sim \mathcal{N}(0, \sigma_A^2)
\end{equation}
where $\sigma_A$ is the standard deviation in dB.

\subsection{Phase Error}
Phase mismatches are modeled as Gaussian deviations from the intended phase:
\begin{equation}
\tilde{\phi}_n \sim \mathcal{N}(\phi_n, \sigma_P^2)
\end{equation}
where $\sigma_P$ is the standard deviation in radians.

\subsection{Delay Error}
TTD control errors are modeled as Gaussian deviations from the intended delay:
\begin{equation}
\tilde{\tau}_n \sim \mathcal{N}(\tau_n, \sigma_T^2)
\end{equation}
where $\sigma_T$ is the standard deviation in picoseconds.
\\ \\
For the baseband TTD implementation the impaired weight vector becomes:
\begin{equation}
[\mathbf{w}_{\text{BB}}[m]]_n = \alpha_n \exp[j(2\pi(f_m - f_c)\tilde{\tau}_n + \tilde{\phi}_n)]
\end{equation}

\section{Results and Discussion}

All simulations were performed using MATLAB with $N_{\text{trials}} = 1000$ Monte Carlo trials. For my implementation The signal-to-noise ratio is defined as $\text{SNR} = \sigma_1^2 / (M \cdot N_R \cdot \sigma_N^2)$, where only the dominant cluster power is used and the noise power is scaled by the number of training subcarriers $M$ and receiver antennas $N_R$.

\subsection{Beam Patterns of TTD Codebook}

Figure~\ref{fig:beam_patterns} shows the beam patterns generated by the designed TTD codebook for $D = 32$ directions with diversity factor $R = 4$. The 32 beams uniformly cover the angular range from $-90°$ to $+90°$ with narrow mainlobes and low sidelobes, demonstrating effective spatial coverage for beam training.

\begin{figure}[H]
\centering
\includegraphics[width=0.85\textwidth]{results/plot1_beam_patterns.png}
\caption{Beam patterns of the designed TTD codebook showing 32 uniformly spaced directional beams.}
\label{fig:beam_patterns}
\end{figure}

\subsection{RMSE vs. SNR}

Figure~\ref{fig:rmse_snr} presents the root mean square error (RMSE) of angle estimation as a function of SNR for diversity factors $R = 1, 2, 4$ without hardware impairments. The diversity factor significantly impacts performance: higher $R$ values provide better estimation accuracy by averaging over more independent channel realizations.

\begin{figure}[H]
\centering
\includegraphics[width=0.75\textwidth]{results/plot2_rmse_vs_snr.png}
\caption{RMSE of angle estimation vs. SNR for diversity factors $R = 1, 2, 4$ without hardware impairments.}
\label{fig:rmse_snr}
\end{figure}

\noindent
At SNR = 0~dB, the RMSE is approximately 15° for $R = 1$, 5° for $R = 2$, and 0.8° for $R = 4$. At high SNR (20~dB), the RMSE approaches but does not reach an error floor determined by the dictionary resolution $\pi/(2Q)$, achieving sub-degree accuracy. The performance improvement from $R = 1$ to $R = 4$ demonstrates the value of frequency diversity in frequency-selective channels.
\\ \\
\textbf{Note:} My results exhibit a higher noise floor compared to Figure 3 in~\cite{boljanovic2020design}, particularly noticeable at high SNR values where my RMSE does reach the error floor. This discrepancy may stem from subtle differences in simulation parameters, random number generation, or numerical precision that are not explicitly detailed in~\cite{boljanovic2020design}, or through some mistake in my implementation. The qualitative trend remains consistent with the reference, but the RMSE values are higher across all SNR ranges.

\subsection{RMSE vs. Gain Error}

Figure~\ref{fig:rmse_gain} illustrates the impact of gain mismatch on angle estimation performance at SNR = 0~dB with diversity factor $R = 4$. The gain error standard deviation $\sigma_A$ ranges from 0~dB to 4.5~dB.

\begin{figure}[H]
\centering
\includegraphics[width=0.75\textwidth]{results/plot3_rmse_vs_gain_error.png}
\caption{RMSE of angle estimation vs. gain error standard deviation $\sigma_A$ at SNR = 0~dB with $R = 4$.}
\label{fig:rmse_gain}
\end{figure}

\noindent
The results show that the algorithm is relatively robust to moderate gain errors. For $\sigma_A < 2$~dB, the RMSE remains below 5°. However, degradation occurs when $\sigma_A > 2$~dB, where the RMSE increases. My gain error results show a qualitatively similar trend to the red dashed line with stars in Figure 4 of~\cite{boljanovic2020design}, but the absolute RMSE values are elevated even when the gain $\sigma_A$ is small due to the noise floor issue. 

\subsection{RMSE vs. Phase Error}

Figure~\ref{fig:rmse_phase} shows the impact of phase mismatch on estimation accuracy at SNR = 0~dB with $R = 4$. The phase error standard deviation $\sigma_P$ ranges from 0° to 50°.

\begin{figure}[H]
\centering
\includegraphics[width=0.75\textwidth]{results/plot4_rmse_vs_phase_error.png}
\caption{RMSE of angle estimation vs. phase error standard deviation $\sigma_P$ at SNR = 0~dB with $R = 4$.}
\label{fig:rmse_phase}
\end{figure}

The algorithm shows good robustness to small phase errors, maintaining RMSE below 5° for $\sigma_P < 20°$. However, performance degrades significantly beyond this threshold, with RMSE exceeding 10° at $\sigma_P = 40°$ and approaching 20° at $\sigma_P = 50°$.

\textbf{Implementation Note:} This result reveals a significant discrepancy with~\cite{boljanovic2020design}. In Figure 4 of the reference paper (red dashed line with diamonds), phase errors have relatively modest impact, with the algorithm maintaining good performance even at $\sigma_P = 50°$. In contrast, my implementation exhibits much stronger sensitivity to phase errors, with rapid degradation beginning around $\sigma_P = 20°$. Furthermore, in the reference paper, gain errors consistently cause more severe degradation than phase errors across all impairment levels. My implementation shows the opposite relationship: phase errors have a substantially larger impact on RMSE than gain errors. This reversal suggests a fundamental difference in how my implementation processes phase information, possibly in the array response calculation, complex exponential evaluation, or the dictionary construction. Despite careful verification of all equations against~\cite{boljanovic2020design}, this behavioral difference remains unexplained and warrants further investigation.

\subsection{RMSE vs. Delay Error}

Figure~\ref{fig:rmse_delay} presents the impact of TTD control errors on baseband (BB) TTD implementation at SNR = 0~dB with $R = 4$. The delay error standard deviation $\sigma_T$ ranges from 0 to 400~ps.

\begin{figure}[H]
\centering
\includegraphics[width=0.75\textwidth]{results/plot5_rmse_vs_delay_error.png}
\caption{RMSE of angle estimation vs. delay error standard deviation $\sigma_T$ for baseband TTD implementation at SNR = 0~dB with $R = 4$. The BB architecture shows excellent robustness to delay errors up to 125~ps.}
\label{fig:rmse_delay}
\end{figure}

The BB TTD architecture exhibits remarkable robustness to delay errors. Performance remains nearly constant for $\sigma_T < 125$~ps, with only gradual degradation up to 200~ps. This resilience is a key advantage of baseband implementation compared to RF TTD arrays, which are highly sensitive to delay errors as reported in~\cite{boljanovic2020design}. The tolerance to delay errors exceeding 100~ps relaxes hardware design constraints significantly, making practical implementation more feasible.

\textbf{Implementation Note:} The delay error results show good qualitative agreement with Figure 5 of~\cite{boljanovic2020design}, exhibiting the expected robustness of BB TTD architecture to delay errors. The absolute RMSE values remain elevated compared to the reference paper, consistent with the noise floor issue observed across all my results. However, the relative behavior and the critical threshold around 125~ps align well with the reference, suggesting that my delay error modeling is consistent with~\cite{boljanovic2020design}.

\subsection{Summary of Findings}

My results demonstrate that:
\begin{enumerate}
    \item The TTD-based beam training algorithm achieves sub-degree to few-degree angle estimation accuracy at moderate to high SNR when $R = 4$.
    \item Frequency diversity significantly improves performance, with each doubling of $R$ providing substantial improvement in estimation accuracy.
    \item The system maintains acceptable performance under various hardware impairments, though with different sensitivities than reported in~\cite{boljanovic2020design}.
    \item Baseband TTD implementation offers substantial advantages in delay error tolerance compared to RF implementation, consistent with the reference paper.
\end{enumerate}

\textbf{Implementation Discrepancies:} Despite faithful implementation of all mathematical formulations from~\cite{boljanovic2020design}, my results exhibit several systematic differences:

\begin{itemize}
    \item \textit{Elevated Noise Floor:} All my results show higher absolute RMSE values compared to the reference paper, particularly noticeable at high SNR. This suggests an unidentified source of excess noise or reduced effective SNR in my simulation, possibly related to numerical precision, random number generation seed effects, or subtle parameter differences not fully specified in~\cite{boljanovic2020design}.

    \item \textit{Reversed Impairment Sensitivity:} The reference paper shows that gain errors cause more severe degradation than phase errors across all impairment levels. My implementation exhibits the opposite behavior, with phase errors having substantially greater impact than gain errors. This reversal is particularly puzzling given that I implemented the impairment models exactly as specified in Equation (5) of~\cite{boljanovic2020design}. Possible explanations include differences in how the dictionary matrix handles phase information, numerical stability issues in complex exponential calculations, or unintended normalization effects in the power-based estimation.

    \item \textit{Qualitative Trends Preserved:} Despite these quantitative differences, the qualitative trends remain consistent: increasing SNR and diversity factor improve performance, and BB TTD architecture provides superior delay error tolerance. This suggests that the core algorithm operates correctly, but some systematic bias affects absolute performance levels and relative sensitivity to different impairment types.
\end{itemize}

These discrepancies highlight the challenges in reproducing simulation results from academic papers, even when mathematical formulations are explicitly provided. Future work should investigate potential sources of these differences, including detailed examination of random number generation, floating-point precision effects, and implicit assumptions in the simulation framework.

\section{Conclusion}

This project successfully implemented and evaluated a single-shot beam training algorithm for wideband millimeter-wave OFDM systems using true-time-delay arrays. Following the methodology in~\cite{boljanovic2020design}, I designed frequency-dependent codebooks, implemented a dictionary-based super-resolution angle estimation algorithm, and comprehensively analyzed performance under realistic hardware impairments.

The key contributions and findings include:
\begin{itemize}
    \item Successful design of TTD codebooks with 32 directional beams covering the full angular range with configurable diversity factor $R = 1, 2, 4$.
    \item Demonstration that frequency diversity significantly improves angle estimation accuracy, achieving few-degree RMSE at high SNR with $R = 4$.
    \item Comprehensive characterization of hardware impairment effects on estimation performance for gain, phase, and delay errors.
    \item Validation that baseband TTD architecture provides superior robustness to delay errors compared to RF implementation, supporting practical deployment.
\end{itemize}

\textbf{Implementation Challenges:} My implementation revealed several systematic discrepancies with the results reported in~\cite{boljanovic2020design}, despite faithful adherence to the mathematical formulations:

\begin{enumerate}
    \item An elevated noise floor across all SNR conditions, resulting in RMSE values 1-2 dB higher than the reference paper.
    \item Reversed relative sensitivity to hardware impairments: phase errors cause more severe degradation in my implementation, whereas the reference paper reports gain errors as more critical.
    \item Nevertheless, qualitative trends remain consistent with~\cite{boljanovic2020design}, and the algorithm demonstrates the expected benefits of frequency diversity and BB TTD architecture.
\end{enumerate}

These observations underscore important lessons about simulation reproducibility in communications research. Even with explicit mathematical models, subtle differences in implementation details—numerical precision, random number generation, implicit parameter assumptions, or normalization procedures—can lead to quantitative discrepancies. The fact that my implementation exhibits correct qualitative behavior but shifted absolute performance suggests that the core algorithm is sound, but some systematic factor affects the noise floor and impairment sensitivity.

Despite these discrepancies, the results confirm that TTD-based beam training offers a viable approach to reducing training overhead in future mmWave and sub-terahertz systems. The single-shot capability represents a significant advantage over conventional exhaustive beam sweeping, and the tolerance to practical hardware impairments (particularly delay errors in BB architecture) supports feasibility for real-world deployment.

Future work should investigate the sources of the observed discrepancies through careful comparison with reference implementations, if available. Additionally, extensions could explore adaptive dictionary refinement, operation at higher carrier frequencies, joint optimization of TTD parameters with hardware constraints, and hybrid analog-digital TTD architectures for improved scalability and energy efficiency. Understanding and resolving the phase error sensitivity in my implementation would be particularly valuable for guiding practical TTD array design.

\begin{thebibliography}{9}

\bibitem{boljanovic2020design}
V.~Boljanovic, H.~Yan, E.~Ghaderi, D.~Heo, S.~Gupta, and D.~Cabric,
``Design of millimeter-wave single-shot beam training for true-time-delay array,''
in \textit{2020 IEEE 21st International Workshop on Signal Processing Advances in Wireless Communications (SPAWC)},
Atlanta, GA, USA, May 2020, pp.~1--5.

\end{thebibliography}

\end{document}